% Generado por Code/tools/tablas_anexos.py desde 8_oe6_qubits.csv. No editar a mano.
\begin{table}[H]
\centering
\small
\caption{Escalamiento computacional con el número de qubits en simulación por vector de estado (OE-6).}
\label{tab:anexo_m8_oe6_qubits}
\begin{adjustbox}{max width=\linewidth}
\begin{tabular}{c c r r r r}
\toprule
$k$ (qubits) & Tipo de dato & Memoria estado & Tiempo/estado (s) & Tiempo embedding (h) & Entrenamiento (h/mapa) \\
\midrule
8 & Medido & 0.00 MiB & 0.0012 & 0.002 & 1.6 \\
10 & Medido & 0.02 MiB & 0.0019 & 0.004 & 3.0 \\
12 & Medido & 0.06 MiB & 0.0029 & 0.006 & 5.7 \\
14 & Medido & 0.25 MiB & 0.0065 & 0.013 & 14.8 \\
16 & Medido & 1.00 MiB & 0.0349 & 0.069 & 90.3 \\
18 & Medido & 4.00 MiB & 0.1071 & 0.212 & 311 \\
20 & Medido & 16.00 MiB & 0.3924 & 0.778 & 1264 \\
22 & Medido & 64.00 MiB & 1.2 & 2.472 & 4415 \\
24 & Medido & 256.00 MiB & 4.1 & 8.194 & 15950 \\
26 & Medido & 1.00 GiB & 13.7 & 27.3 & 57421 \\
28 & Medido & 4.00 GiB & 74.8 & 148.2 & 336114 \\
30 & Extrapolado & 16.00 GiB & 221.0 & 438.0 & 1063514 \\
32 & Extrapolado & 64.00 GiB & 802.7 & 1591.1 & 4118952 \\
\bottomrule
\end{tabular}
\end{adjustbox}
\par\smallskip
\parbox{0.95\linewidth}{\footnotesize Escalamiento del coste en simulación cuántica por vector de estado en función del número de qubits ($k$). Para $k \le 28$ las cifras corresponden a ejecuciones empíricas en el hardware de referencia (Apple M3); para $k \ge 30$ corresponden a la extrapolación exponencial teórica.}
\end{table}
